\documentclass[12pt]{report}
\usepackage[T1]{fontenc}
\usepackage[brazilian]{babel}
\usepackage{csquotes}
\usepackage{makeidx}
\usepackage{glossaries}
\usepackage[
  perfil=academico,
  tipo-trabalho=tcc,
  impressao=frente-e-verso,
  citacoes=autor-data
]{abntex3}

\makeindex
\makeglossaries
\newglossaryentry{software}{
  name={software livre},
  description={programa que assegura liberdades de uso, estudo e modificação}
}
\addbibresource{referencias-exemplo.bib}
\ABNTEXsetup{
  titulo={Modelo canônico completo de trabalho de conclusão de curso},
  subtitulo={base auditável com todos os elementos implementados},
  autoria={Maria da Silva},
  instituicao={Universidade Exemplo},
  natureza={Trabalho de conclusão de curso},
  objetivo={apresentado para obtenção do grau de Bacharela},
  area={Sistemas de Informação},
  orientacao={Prof. Dr. João Souza},
  coorientacao={Profa. Dra. Ana Santos},
  local={Fortaleza},
  data={2026}
}
\ABNTEXabreviatura{cap.}{capítulo}
\ABNTEXsigla{TCC}{Trabalho de Conclusão de Curso}
\ABNTEXsimbolo{$\%$}{Percentual}

\begin{document}
\ABNTEXacademicoexterna
\ABNTEXlombada[TCC 2026][Universidade Exemplo]

\ABNTEXacademicopretextual
\ABNTEXfolhaderosto
\ABNTEXfichacatalografica{%
  Ponto de extensão para dados catalográficos, quando exigidos.
}
\ABNTEXerrata
  {SILVA, Maria da. \textbf{Modelo canônico completo de trabalho de
   conclusão de curso}. 2026. TCC — Universidade Exemplo, Fortaleza, 2026.}
  {\begin{center}
     \begin{tabular}{llll}
       Folha & Linha & Onde se lê & Leia-se\\
       8 & 2 & sofware & software
     \end{tabular}
   \end{center}}
\ABNTEXfolhadeaprovacao
  {30 de julho de 2026}
  {%
    \ABNTEXmembrobanca
      {Prof. Dr. João Souza}{Orientador}{Universidade Exemplo}
    \ABNTEXmembrobanca
      {Profa. Dra. Ana Santos}{Coorientadora}{Universidade Exemplo}
    \ABNTEXmembrobanca
      {Prof. Me. Carlos Lima}{Examinador}{Instituto Exemplo}
  }
\ABNTEXdedicatoria{À família e à comunidade acadêmica.}
\ABNTEXagradecimentos{À orientação, à banca e às pessoas colaboradoras.}
\ABNTEXepigrafe[Autoria demonstrativa]{Testar também é documentar.}
\ABNTEXresumo
  {Modelo completo de TCC que exercita os elementos públicos implementados e
   fornece pontos objetivos para auditoria externa.}
  {TCC, modelo, auditoria, LaTeX}
\ABNTEXresumoemlinguaestrangeira
  {Complete undergraduate final project template exercising implemented
   public elements and providing objective external audit points.}
  {final project, template, audit, LaTeX}
\ABNTEXlistadeilustracoes
\ABNTEXlistadetabelas
\ABNTEXlistadeabreviaturasesiglas
\ABNTEXlistadesimbolos
\ABNTEXsumario

\ABNTEXacademicotextual
\chapter{Introdução}
\index{software}
O capítulo apresenta tema, problema, objetivos e justificativa. O
\gls{software} é usado apenas como termo de demonstração
\parencite{associacao2025}.
\section{Problema}
Como produzir documentos acadêmicos verificáveis com uma API pública?
\section{Objetivos}
Demonstrar elementos e registrar evidências de compilação.

\chapter{Fundamentação teórica}
\textcite{associacao2025} demonstra a citação narrativa.
\begin{ABNTEXcitacaolonga}
Texto demonstrativo destinado a permitir inspeção visual do ambiente de
citação longa sem incorporar material protegido.
\end{ABNTEXcitacaolonga}
\ABNTEXnota{Nota demonstrativa do modelo de TCC.}

\chapter{Metodologia}
\section{Etapas}
Inventário, composição, compilação e inspeção constituem as etapas.

\chapter{Desenvolvimento}
\section{Implementação}
O documento emprega somente comandos públicos.
\begin{figure}[ht]
  \centering
  \rule{0.55\linewidth}{2cm}
  \caption{Protótipo demonstrativo}
  \ABNTEXfonte{elaborada pela autora.}
\end{figure}
\begin{table}[ht]
  \centering
  \caption{Etapas de validação}
  \begin{tabular}{ll}
    Etapa & Situação\\
    Fonte & versionada\\
    PDF & compilado
  \end{tabular}
  \ABNTEXfonte{elaborada pela autora.}
\end{table}

\chapter{Resultados e discussão}
Registram-se resultados, confronto com a literatura e limitações.

\chapter{Conclusão}
Sintetizam-se as respostas ao problema e os trabalhos futuros.

\ABNTEXacademicopostextual
\ABNTEXbibliografia
\ABNTEXglossario
\ABNTEXappendix{Roteiro de testes elaborado pela autora}
Conteúdo produzido no desenvolvimento do TCC.
\ABNTEXannex{Especificação externa de demonstração}
Conteúdo de origem externa.
\ABNTEXindice[Índice de assuntos]
\end{document}
